\documentclass[12pt]{article}
\usepackage[utf8]{inputenc}
\usepackage[T1]{fontenc}
\usepackage{hyperref}
\title{Murut, Paluan Language Resources for K-12 Curriculum Development}
\author{Compiled by Claude Sonnet 5.5 (claude-sonnet-5-5) for Andr\'as Kornai}
\date{}
\begin{document}
\maketitle

\section{Introduction}
Paluan Murut (ISO 639-3 \texttt{plz}, Glottocode palu1253) is a Murutic Austronesian language of Sabah, Malaysia. Its main dialects are Paluan (Peluan) and Pandewan. About 5,500 speakers are reported. Few dedicated resources were found.

\section{Primary Educational Resources}
No primers, textbooks or school programmes for Paluan were found.

\section{Linguistic and Reference Materials}
\begin{itemize}
\item Wikipedia article on the Paluan language: \url{https://en.wikipedia.org/wiki/Paluan_language}
\item Wikipedia article on the Murutic languages: \url{https://en.wikipedia.org/wiki/Murutic_languages}
\item ASJP wordlist for Murut Paluan, from Lobel (2016), \emph{North Borneo Sourcebook}: \url{https://asjp.clld.org/languages/MURUT_PALUAN}
\item Austronesian Comparative Dictionary entry for Paluan: \url{https://acd.clld.org/languages/917}
\item Prentice (1971), \emph{The Murut languages of Sabah}, listed in the ACD sources. Its coverage of Paluan was not confirmed: \url{https://acd.clld.org/sources/prentice1971}
\end{itemize}

\section{Religious and Community Resources}
\begin{itemize}
\item Joshua Project profile of Paluan Murut in Malaysia (demographic and religious information only): \url{https://joshuaproject.net/people_groups/14354}
\end{itemize}

\section{Digital and Audio Resources}
No Paluan audio or software resources were found. A related but different language, Timugon Murut, has an online Webonary dictionary of about 22,600 words, which shows how a Paluan lexicon could be hosted:
\begin{itemize}
\item \url{https://www.webonary.org/timugon_murut/en/overview/introduction/}
\item \url{https://fmtv5.freemalaysiatoday.com/category/nation/2026/01/02/language-activists-put-timugon-murut-dictionary-online}
\end{itemize}

\section{Curriculum Development Context}
Paluan is not known to be taught in schools. A curriculum would have to start with documentation (a dictionary and orthography). The Timugon project is a possible model, and its materials must not be mistaken for Paluan.
\end{document}